\section{Experiments}

\subsection{Protocol}

We use the optimized Kalman filter as the primary baseline and compare it against residual models. Unless otherwise noted, models use a $400$ ms history window, an $800$ ms forecast horizon, normalized boxes, AdamW training with learning rate $10^{-3}$, weight decay $10^{-4}$, Smooth L1 loss, and gradient clipping at norm $10$. Development runs use a folder-level train-evaluation split with $20\%$ of the training folders held out. The final held-out test split is evaluated only after model selection.

The main image experiments use fixed cutouts, relative box-history features, frame-centered filtered center positions, optimized Kalman state features, and ResNet-18 encoders. We do not use Kalman covariance features in the main results. Because training time is limited, we report a compact set of representation comparisons rather than a full factorial sweep.

\subsection{Kalman and Linear Diagnostics}

The first stage establishes how much performance is available without images. We compare the last-four linear extrapolator, the optimized center-size Kalman filter, and linear residual heads over selected filtered state features. These linear residual models use no image inputs and no box-history MLP.

\begin{table}[t]
  \centering
  \caption{Non-image forecasting baselines and linear residual diagnostics. Lower is better for ADE/FDE; higher is better for mIoU.}
  \label{tab:linear-diagnostics}
  \resizebox{\textwidth}{!}{%
  \begin{tabular}{lcccccc}
    \toprule
    Method & Learned features & Size residuals & ADE-C $\downarrow$ & FDE-C $\downarrow$ & ADE-BB $\downarrow$ & mIoU $\uparrow$ \\
    \midrule
    Last-four CV & none & no & \todo{} & \todo{} & \todo{} & \todo{} \\
    Optimized Kalman & none & no & \todo{} & \todo{} & \todo{} & \todo{} \\
    Linear residual & center position & yes/no & \todo{} & \todo{} & \todo{} & \todo{} \\
    Linear residual & center velocity & yes/no & \todo{} & \todo{} & \todo{} & \todo{} \\
    Linear residual & center position+velocity & yes/no & \todo{} & \todo{} & \todo{} & \todo{} \\
    Linear residual & velocities & yes/no & \todo{} & \todo{} & \todo{} & \todo{} \\
    \bottomrule
  \end{tabular}
  }
\end{table}

Preliminary runs indicate that filtered position and velocity can still predict useful acceleration residuals, even after decorrelation, although the resulting center improvements do not always translate into better mIoU. We therefore treat these models as diagnostics for motion-prior leakage rather than as the main contribution.

\subsection{Motion-Prior Correlation}

We compute acceleration-field summaries on the training, train-evaluation, and held-out test splits. For each split, we report mean absolute Pearson correlation and ridge-linear $R^2$ from motion-prior features to fitted center acceleration. We also visualize acceleration as a function of image position and velocity bins to expose edge-to-center trends, speed-dependent deceleration, and global vertical bias.

\begin{table}[t]
  \centering
  \caption{Motion-prior predictability of fitted future acceleration before and after sample decorrelation.}
  \label{tab:correlation}
  \begin{tabular}{lccccc}
    \toprule
    Split & Decorr. & Samples & Mean $|\rho|$ $\downarrow$ & Mean $R^2$ $\downarrow$ & $\|\bar{\mathbf{a}}\|_2$ $\downarrow$ \\
    \midrule
    Train & no & \todo{} & \todo{} & \todo{} & \todo{} \\
    Train & yes & \todo{} & \todo{} & \todo{} & \todo{} \\
    Train-eval & no & \todo{} & \todo{} & \todo{} & \todo{} \\
    Train-eval & yes & \todo{} & \todo{} & \todo{} & \todo{} \\
    Test & no & \todo{} & \todo{} & \todo{} & \todo{} \\
    Test & yes & \todo{} & \todo{} & \todo{} & \todo{} \\
    \bottomrule
  \end{tabular}
\end{table}

\subsection{Standard-Split Image Results}

The second stage evaluates event-conditioned residuals on the standard split. Early experiments suggest that event-derived inputs improve box metrics beyond the optimized Kalman baseline on this split. Dataset-native event frames have been the strongest single image input in the raw split among the tested formats, while CSTR2, CSTR3, XT, and YT provide complementary views of motion. We have not yet completed a broad search over all combinations, so the final table will focus on the strongest single inputs and a small number of combinations.

\begin{table}[t]
  \centering
  \caption{Standard-split forecasting results before decorrelation. Lower is better for ADE/FDE; higher is better for mIoU.}
  \label{tab:standard-results}
  \resizebox{\textwidth}{!}{%
  \begin{tabular}{lcccccc}
    \toprule
    Method & Inputs & Cutout & ADE-C $\downarrow$ & FDE-C $\downarrow$ & ADE-BB $\downarrow$ & mIoU $\uparrow$ \\
    \midrule
    Optimized Kalman & boxes & none & \todo{} & \todo{} & \todo{} & \todo{} \\
    Residual & filter state only & none & \todo{} & \todo{} & \todo{} & \todo{} \\
    Residual & event frames & fixed & \todo{} & \todo{} & \todo{} & \todo{} \\
    Residual & CSTR2 & fixed & \todo{} & \todo{} & \todo{} & \todo{} \\
    Residual & CSTR3 & fixed & \todo{} & \todo{} & \todo{} & \todo{} \\
    Residual & XT+YT & fixed & \todo{} & \todo{} & \todo{} & \todo{} \\
    Residual & event/CSTR + XT/YT & fixed & \todo{} & \todo{} & \todo{} & \todo{} \\
    Residual & event/CSTR + XT/YT + RGB & fixed & \todo{} & \todo{} & \todo{} & \todo{} \\
    \bottomrule
  \end{tabular}
  }
\end{table}

\subsection{Decorrelated-Subset Results}

The third stage repeats the strongest standard-split settings on the decorrelated sample pools. The decorrelation procedure keeps $50\%$ of capped samples, uses the motion-prior feature mode, and applies a small mean-acceleration penalty. Preliminary results suggest that decorrelation reduces but does not eliminate residual predictability from position and velocity. On decorrelated samples, CSTR2 has so far appeared stronger than dataset-native event frames, possibly because its normalized timestamp structure exposes local motion gradients that binary-like frames do not. This explanation remains speculative and should be treated as a hypothesis for follow-up experiments.

\begin{table}[t]
  \centering
  \caption{Forecasting results on the decorrelated sample split.}
  \label{tab:decorrelated-results}
  \resizebox{\textwidth}{!}{%
  \begin{tabular}{lcccccc}
    \toprule
    Method & Inputs & Keep frac. & ADE-C $\downarrow$ & FDE-C $\downarrow$ & ADE-BB $\downarrow$ & mIoU $\uparrow$ \\
    \midrule
    Optimized Kalman & boxes & 0.5 & \todo{} & \todo{} & \todo{} & \todo{} \\
    Residual & filter state only & 0.5 & \todo{} & \todo{} & \todo{} & \todo{} \\
    Residual & event frames & 0.5 & \todo{} & \todo{} & \todo{} & \todo{} \\
    Residual & CSTR2 & 0.5 & \todo{} & \todo{} & \todo{} & \todo{} \\
    Residual & CSTR2 + XT/YT & 0.5 & \todo{} & \todo{} & \todo{} & \todo{} \\
    Residual & best available combo & 0.5 & \todo{} & \todo{} & \todo{} & \todo{} \\
    \bottomrule
  \end{tabular}
  }
\end{table}

\subsection{Comparison to Prior FRED Results}

We compare against the original FRED box metrics and the recent center-forecasting results from \todo{cite prior arXiv work}. Our optimized Kalman extension to width and height appears to improve over the original FRED box results on mIoU and box ADE/FDE, while center ADE/FDE remains competitive but does not clearly surpass the strongest reported center-only methods. \todo{Insert exact numbers and ensure metric definitions match prior work.}

\subsection{Ablations Deferred by Scope}

The following ablations are useful but not required for the core workshop claim: cutout size, temporal bin count, temporal crop length for XT/YT, smaller CNN backbones, covariance features, RGB contribution, and broader representation combinations. We report them as limitations or future work unless completed before submission.
